\chapter{Geometric semistability and nonzero pairings}

Let $K$ be a field of characteristic zero, let $J$ be a finite set, and let
$(E_j)_{j\in J}$ and $V$ be nonzero finite-dimensional $K$-vector spaces.
Write
\[
 W=\bigotimes_{j\in J}E_j,\qquad a_j=\dim_K E_j,\qquad s=\dim_K V,
 \qquad M^\vee=\Hom_K(M,K),\qquad \det M=\bigwedge^{\dim_K M}M.
\]
All tensor products in this chapter are algebraic. If $L/K$ is an extension,
put $M_L=L\otimes_K M$ and use the canonical maps
\[
 \beta_E:L\otimes_K W\xrightarrow{\sim}\bigotimes_j(E_j)_L,
 \qquad
 \delta_M:L\otimes_K M^\vee\xrightarrow{\sim}(M_L)^\vee.
\]
They are characterized by
$\beta_E(1\otimes\bigotimes_j x_j)=\bigotimes_j(1\otimes x_j)$ and
$\delta_M(c\otimes\ell)(b\otimes x)=cb\ell(x)$.
For $A:V\to W$ write $A_L=\beta_E(1\otimes A)$; for $q:W\to Q$ write
$q_L=(1\otimes q)\beta_E^{-1}$.
We develop the invariant-polynomial argument in the proof of Theorem 6.2 of
this paper \cite{bost2026}, keeping these maps fixed throughout.

\section{Polynomial and Pl\"ucker models}

\begin{lemma}[Coordinate and intrinsic invariant polynomials]
\label{lem:invariant-polynomial-models}
\lean{BostConjecture.Adelic.Geometry.isPolynomialSemistable_iff_exists_homogeneousInvariant}
Let $F$ be an infinite field, $X$ a finite-dimensional $F$-vector space,
$b=(b_\alpha)_{\alpha\in I}$ a basis, and $\rho:G\to\GL(X)$ a group action.
A vector $x\in X$ admits a positive-degree homogeneous coordinate polynomial
$P\in F[T_\alpha\mid\alpha\in I]$ such that
\[
 P(b^{-1}\rho(g)y)=P(b^{-1}y)\quad(g\in G,\ y\in X),
 \qquad P(b^{-1}x)\ne0,
\]
if and only if there is a positive-degree homogeneous
$p\in\Sym_F(X^\vee)$ fixed by pullback under every $\rho(g)$ with $p(x)\ne0$.
\end{lemma}
\begin{proof}
Send the dual basis vector $b_\alpha^\vee$ to $T_\alpha$. The resulting
algebra isomorphism $\Sym_F(X^\vee)\simeq F[T_\alpha]$ preserves degree and
sends evaluation at $y$ to evaluation at $b^{-1}y$. Pullback by a linear map
$f$ is substitution $\ell\mapsto\ell\circ f$, so its evaluation at $y$ is
$p(fy)$. Equality of polynomials follows from equality of their values over
an infinite field: in one variable a nonzero polynomial has at most its
degree many roots; in several variables, expand in the last variable and
apply that fact, then induct on the number of variables to each coefficient.
Thus the displayed pointwise invariance is exactly equality under pullback.
The same isomorphism preserves the specified nonzero value, proving both
implications.
\end{proof}

\begin{definition}[Geometric semistability]
\label{def:geometric-semistability}
\source{bost-adelic-2026:page-0026,bost-adelic-2026:page-0028}
\uses{lem:invariant-polynomial-models}
\lean{ArakelovGeometry.InvariantTheory.IsGeometricallySemistableTensorSubspace,ArakelovGeometry.InvariantTheory.IsGeometricallySemistableTensorQuotient}
Fix an algebraic closure $\overline K$ and put
$G=\prod_j\SL((E_j)_{\overline K})$, acting on
$W_{\overline K}=\bigotimes_j(E_j)_{\overline K}$ by tensoring its factors.
A map $A:V\to W$ is \emph{geometrically semistable} if there exist
$0\ne\eta\in\det(V_{\overline K})$ and a homogeneous polynomial of degree
$m>0$ on $\bigwedge^s W_{\overline K}$, invariant under the induced $G$-action,
whose value at $(\bigwedge^s A_{\overline K})(\eta)$ is nonzero.
For a subspace this definition is applied to its inclusion.
For a nonzero finite-dimensional $Q$, a map $q:W\to Q$ is geometrically
semistable if $q_{\overline K}^\vee:Q_{\overline K}^\vee\to W_{\overline K}^\vee$
has the same Pl\"ucker property under the contragredient action
$g\ell=\ell\circ g^{-1}$. For a quotient this definition is applied to its
surjection.
\end{definition}

The quantifier over $G$ here includes all special linear transformations over
$\overline K$. Homogeneity gives $p(cz)=c^m p(z)$, so the condition is
unchanged upon replacing a determinant generator by a nonzero multiple.
For $s>0$ it forces $A$ to be injective: a noninjective map has zero $s$th
exterior power, while a positive-degree homogeneous polynomial vanishes at
zero. Applied to the dual, this also forces $q$ to be surjective.

\begin{lemma}[Pl\"ucker substitution]
\label{lem:plucker-invariant-substitution}
\source{bost-adelic-2026:page-0027}
\uses{def:geometric-semistability,lem:invariant-polynomial-models}
\lean{BostConjecture.Adelic.Geometry.exists_homogeneousInvariant_of_isPluckerSemistable,BostConjecture.Geometry.pluckerPolynomialPullback_mem_invariants,BostConjecture.Geometry.polynomialEval_pluckerPolynomialPullback_ne_zero}
Let $F$ be an infinite field, $X$ a finite-dimensional $F$-vector space with
an action of $G$, and $V$ an $s$-dimensional $F$-vector space.
For a basis $(v_1,\ldots,v_s)$ of $V$ put $\eta=v_1\wedge\cdots\wedge v_s$.
If $p$ is a homogeneous degree-$m$ invariant on $\bigwedge^sX$, then
\[
 Q_p(B)=p\bigl(Bv_1\wedge\cdots\wedge Bv_s\bigr),\qquad B\in\Hom_F(V,X),
\]
is a homogeneous degree-$sm$ invariant for
$(g,h)B=\rho(g)Bh^{-1}$, where $(g,h)\in G\times\SL(V)$.
If $A$ is Pl\"ucker semistable, $m>0$ and $p$ can be chosen with $Q_p(A)\ne0$
for this prescribed basis.
\end{lemma}
\begin{proof}
Each coordinate of $Bv_1\wedge\cdots\wedge Bv_s$ is a determinant of an
$s$-by-$s$ matrix of linear coordinates of $B$, hence is homogeneous of
degree $s$. Substitution into $p$ gives degree $sm$.
Since $\bigwedge^s h^{-1}$ acts on $\det V$ by $\det(h)^{-1}=1$,
\[
 (\bigwedge^s(\rho(g)Bh^{-1}))(\eta)
   =(\bigwedge^s\rho(g))(\bigwedge^s B)(\eta).
\]
Invariance of $p$ proves invariance of $Q_p$. For the last assertion, write
the nonzero determinant generator supplied by semistability as $c\eta$,
where $c\ne0$. Its nonzero polynomial value is
$c^m p((\bigwedge^s A)(\eta))$, so the latter value is nonzero as well.
The coordinate-to-intrinsic identification preserves this value.
\end{proof}

\section{The original-field pairing}

For $a,k>0$ and $d=ak$, a rectangular tableau $t$ assigns the tensor
positions $\{1,\ldots,d\}$ to $a$ rows and $k$ columns. In
$(\mathbb Z^a)^{\otimes d}$ let $w_t$ be the product, in these positions,
of $k$ copies of the alternating tensor
\[
 \sum_{\sigma\in\mathfrak S_a}\operatorname{sgn}(\sigma)
 e_{\sigma(1)}\otimes\cdots\otimes e_{\sigma(a)}.
\]
Let $L_t$ be the $\mathbb Z$-span of its position-permutation orbit.
For an $a$-dimensional space $M$ set
\[
 S_t(M)=(\det M)^{\otimes k}\otimes_K(K\otimes_{\mathbb Z}L_t).
\]
The determinant map $\iota_{M,t}:S_t(M)\to M^{\otimes d}$ sends
$(b_1\wedge\cdots\wedge b_a)^{\otimes k}\otimes x$ to the tensor obtained
by substituting $b_i$ for every letter $e_i$ of $x$.
Lemma~\ref{lem:rectangular-specht-invariants} identifies its image and ensures
that the construction is independent of the chosen basis.
There is no division by $a!$ in this definition.

\begin{definition}[Canonical determinant pairing]
\label{def:canonical-pairing}
\source{bost-adelic-2026:page-0025,bost-adelic-2026:page-0026}
\uses{lem:rectangular-specht-invariants}
\lean{BostConjecture.Geometry.canonicalPairing}
Let $d>0$ be divisible by $s$ and every $a_j$, put $l=d/s$ and $k_j=d/a_j$,
and choose rectangular tableaux $v$ of shape $(l^s)$ and $t_j$ of shape
$(k_j^{a_j})$. Define
\[
 B_d=(\det V)^{\otimes l}\otimes_K(K\otimes_{\mathbb Z}L_v),
 \qquad
 D_d=\bigotimes_j\left((\det E_j^\vee)^{\otimes k_j}
              \otimes_K(K\otimes_{\mathbb Z}L_{t_j})\right).
\]
Let $j_d:D_d\to(W^\vee)^{\otimes d}$ be the tensor product of the maps
$\iota_{E_j^\vee,t_j}$ followed by regrouping and the tensor-duality map
\[
 e_E:\bigotimes_j E_j^\vee\xrightarrow{\sim}W^\vee,
 \qquad e_E(\bigotimes_j\ell_j)(\bigotimes_jx_j)=\prod_j\ell_j(x_j).
\]
The \emph{canonical pairing of $A$} is the $K$-linear map
\[
 \Psi_{A,d}:B_d\longrightarrow D_d^\vee,\qquad
 \Psi_{A,d}(z)(u)=
 \left\langle j_d(u),A^{\otimes d}\iota_{V,v}(z)\right\rangle.
\]
Evaluation on pure tensors is
$\langle\ell_1\otimes\cdots\otimes\ell_d,
 x_1\otimes\cdots\otimes x_d\rangle=\prod_r\ell_r(x_r)$.
Thus $\Psi_{A,d}=j_d^\vee\circ A^{\otimes d}\circ\iota_{V,v}$ under the
canonical tensor-duality identification.
\end{definition}

\begin{lemma}[Nonzero pairings in an unbounded sequence of degrees]
\label{lem:nonzero-pairings}
\source{bost-adelic-2026:page-0027}
\uses{def:geometric-semistability,lem:plucker-invariant-substitution,def:canonical-pairing,lem:rectangular-specht-invariants}
\lean{BostConjecture.Adelic.Geometry.exists_canonicalPairing_nonzero_sequence_of_geometricSemistability}
If $A:V\to W$ is geometrically semistable, there exists an integer $d_0>0$
divisible by $s$ and by every $a_j$, such that for every integer $r>0$ there
are tableaux as in Definition~\ref{def:canonical-pairing} at degree $d=d_0r$
for which the original $K$-linear map $\Psi_{A,d}$ is nonzero.
\end{lemma}
\begin{proof}
Work temporarily over $F=\overline K$. Choose a basis of $V_F$ and use
Lemma~\ref{lem:plucker-invariant-substitution} to obtain $m>0$ and a
homogeneous invariant $Q$ of degree $sm$ on $\Hom_F(V_F,W_F)$, with
$Q(A_F)\ne0$. Put $P=\prod_j a_j$ and $d_0=smP$. Then $d_0>0$, $s\mid d_0$,
and $a_j\mid d_0$ for every $j$. For $r>0$ the polynomial $Q^{Pr}$ is
homogeneous of degree $d=d_0r$, is invariant, and is nonzero at $A_F$.
This choice enlarges the degree used on page 27 of this paper; it ensures
all divisibilities directly, without a primitive-root argument.

Here is the tensor representative and its invariance in detail. For any
finite-dimensional space $H$, the quotient map
$\pi:(H^\vee)^{\otimes d}\to\Sym^d(H^\vee)$ has a section
\[
 \operatorname{sym}(\ell_1\cdots\ell_d)
 =\frac1{d!}\sum_{\sigma\in\mathfrak S_d}
       \ell_{\sigma(1)}\otimes\cdots\otimes\ell_{\sigma(d)}.
\]
The formula is multilinear and invariant under permutation of the factors,
so it descends through the symmetric quotient; $d!$ is invertible in $F$.
It satisfies $\pi\operatorname{sym}=\id$ and
$\operatorname{sym}\pi=(d!)^{-1}\sum_\sigma\sigma$.
Consequently $\pi$ restricts to an isomorphism from the position-invariant
tensors to $\Sym^d(H^\vee)$. Tensoring a linear map commutes with every
position permutation. Hence this section intertwines every linear group
action, and an invariant polynomial has an invariant tensor representative.
Evaluation of the representative on $B^{\otimes d}$ is its polynomial value
at $B$, because each summand of the average has that same value.

Apply this with $H=\Hom_F(V_F,W_F)$ and regroup
$H^\vee=W_F^\vee\otimes_F V_F$. The representative of $Q^{Pr}$ belongs to
\[
 \left(\bigotimes_j((E_j)_F^\vee)^{\otimes d}
                 \otimes_F V_F^{\otimes d}\right)
 ^{\left(\prod_j\SL((E_j)_F)\right)\times\SL(V_F)}.
\]
Invariants under independent groups can be taken separately. Indeed, if a
group acts only on $X$ and $y_1,\ldots,y_n$ is a basis of $Y$, then
$\sum_i x_i\otimes y_i$ is fixed precisely when each $x_i$ is fixed, as
seen by applying the dual basis of $Y$. Repeat this argument factor by
factor. The contragredient maps $g\mapsto(g^{-1})^\vee$ exhaust
$\SL((E_j)_F^\vee)$: transposition and inverse preserve determinant one
and are bijective. Lemma~\ref{lem:rectangular-specht-invariants} therefore
identifies these independent invariant spaces with the images of
$(D_d)_F$ and $(B_d)_F$. Rectangular tableaux exist since
$d=a_jk_j=sl$ with $k_j,l>0$. We may thus write the representative as
$\sum_\nu j_{d,F}(u_\nu)\otimes\iota_{V_F,v}(z_\nu)$, with
$u_\nu\in(D_d)_F$ and $z_\nu\in(B_d)_F$. Evaluating gives
\[
 0\ne Q(A_F)^{Pr}=\sum_\nu\Psi_{A_F,d}(z_\nu)(u_\nu).
\]

Finally the displayed pairing over $F$ is the scalar extension of the one
over $K$. For the determinant maps this follows by writing tensors in a
$K$-basis: substitution of the basis letters and the integral alternating
coefficients commutes with scalar extension. Regrouping commutes with scalar
extension on pure tensors. The duality maps commute because their evaluations
are products of the same scalars. Pure tensors span every tensor product,
so these equalities prove compatibility on the entire source and target.
If $\Psi_{A,d}$ were zero, its scalar extension would be zero, contradicting
the last display. This proves the assertion over the original field.
\end{proof}

\begin{lemma}[Simultaneous permutation equivariance]
\label{lem:canonical-pairing-equivariance}
\source{bost-adelic-2026:page-0027}
\uses{def:canonical-pairing}
\lean{BostConjecture.Geometry.canonicalPairing_reindex_left}
For $\sigma\in\mathfrak S_d$, let $\sigma$ act on $B_d$ through its action
on $L_v$ and on $D_d$ through its simultaneous actions on all $L_{t_j}$,
leaving determinant lines fixed. Then
\[
 \Psi_{A,d}(\sigma z)(u)=\Psi_{A,d}(z)(\sigma^{-1}u).
\]
\end{lemma}
\begin{proof}
The determinant maps send a position permutation of an orbit-lattice tensor
to the same position permutation of its image. Regrouping sends simultaneous
permutations in all factors to that position permutation in
$(W^\vee)^{\otimes d}$. The map $A^{\otimes d}$ commutes with these
permutations. On pure tensors, reindexing the finite product defining
evaluation gives $\langle u,\sigma x\rangle=\langle\sigma^{-1}u,x\rangle$.
Bilinearity extends this identity to all tensors, proving the formula.
\end{proof}

\section{Quotients and duality}

\begin{lemma}[Scalar extension of the dual inclusion]
\label{lem:dual-inclusion-base-change}
\lean{BostConjecture.Adelic.Geometry.dualTensorInclusion_baseChange}
For any field extension $L/K$ and linear map $q:W\to Q$, define
$A_q=e_E^{-1}q^\vee:Q^\vee\to\bigotimes_jE_j^\vee$.
With the canonical maps fixed above, one has
\[
 \left(\bigotimes_j\delta_{E_j}\right)\beta_{E^\vee}(1\otimes A_q)
   =e_{E_L}^{-1}q_L^\vee\delta_Q.
\]
\end{lemma}
\begin{proof}
Compose both sides with $e_{E_L}$ and evaluate on pure tensors
$\bigotimes_j(b_j\otimes x_j)$. For an input $c\otimes\ell$ both sides
have value
$c(\prod_j b_j)\ell(q(\bigotimes_jx_j))$ by the defining formulas of
$\beta$, $\delta$ and $e$. Pure tensors span the source and the space on
which the resulting functionals are evaluated. The two maps after composition
with $e_{E_L}$ are therefore equal; its injectivity proves the assertion.
\end{proof}

\begin{lemma}[Geometric semistability of the original dual inclusion]
\label{lem:geometric-dual-quotient}
\source{bost-adelic-2026:page-0028}
\uses{def:geometric-semistability,lem:dual-inclusion-base-change}
\lean{BostConjecture.Adelic.Geometry.geometricSemistability_dualTensor_of_quotient}
If $q:W\to Q$ is geometrically semistable, then its original-field dual
inclusion
\[
 A_q=e_E^{-1}q^\vee:Q^\vee\longrightarrow\bigotimes_jE_j^\vee
\]
is geometrically semistable for the full product of special linear groups
of the dual factors.
\end{lemma}
\begin{proof}
Put $F=\overline K$. The map $e_{E_F}$ intertwines the action
$\bigotimes_j(g_j^{-1})^\vee$ on $\bigotimes_j(E_j)_F^\vee$ with the
contragredient tensor action on $W_F^\vee$: evaluation on
$\bigotimes_jx_j$ gives $\prod_j\ell_j(g_j^{-1}x_j)$ on either side.
Every tuple of special linear maps on the dual factors has this form,
since $g\mapsto(g^{-1})^\vee$ is bijective. Pull back the invariant
polynomial witnessing semistability of $q_F^\vee$ along
$\bigwedge^{\dim Q}e_{E_F}$. Its degree and invariant property are preserved,
and its value at the determinant point of $e_{E_F}^{-1}q_F^\vee$ is the
original nonzero value.

The isomorphisms $\delta_{E_j}$ identify $F\otimes_KE_j^\vee$ with
$(E_j)_F^\vee$ and conjugate their full special linear groups: conjugation
preserves determinant and is bijective. Lemma~\ref{lem:dual-inclusion-base-change}
identifies the map just tested with the scalar extension of $A_q$, up to
these target isomorphisms and the source isomorphism $\delta_Q$.
Exterior powers carry the corresponding determinant point to the same point.
Pulling back the polynomial along the target isomorphism and the nonzero
determinant generator along the source isomorphism proves the asserted
geometric semistability of $A_q$ itself.
\end{proof}

\begin{lemma}[Binary tensor-duality convention]
\label{lem:binary-dual-inclusion}
\lean{BostConjecture.Adelic.Geometry.dualTensorInclusion_bool}
For $J=\{0,1\}$, let $b_E:\bigotimes_{j\in J}E_j\to E_0\otimes_KE_1$
be the equivalence sending a pure tensor to $x_0\otimes x_1$; define
$b_{E^\vee}$ in the same way. If $q:E_0\otimes_KE_1\to Q$, then
\[
 b_{E^\vee}\, e_E^{-1}(q b_E)^\vee
       = e_{0,1}^{-1}q^\vee,
\]
where $e_{0,1}:E_0^\vee\otimes_KE_1^\vee\to(E_0\otimes_KE_1)^\vee$
is the ordinary binary tensor-duality map.
\end{lemma}
\begin{proof}
For $\ell_0,\ell_1,x_0,x_1$, evaluation through either tensor convention is
$\ell_0(x_0)\ell_1(x_1)$. Thus
$e_{0,1}(b_{E^\vee}z)(b_Ex)=e_E(z)(x)$ for all $z,x$, by bilinearity
and spanning by pure tensors. Taking $z=e_E^{-1}(qb_E)^\vee\ell$ gives
$e_{0,1}(b_{E^\vee}z)(b_Ex)=\ell(q(b_Ex))$.
Surjectivity of $b_E$ and injectivity of $e_{0,1}$ prove the identity.
\end{proof}

\begin{corollary}[Nonzero pairings for a geometric quotient]
\label{cor:geometric-quotient-nonzero-pairings}
\source{bost-adelic-2026:page-0027,bost-adelic-2026:page-0028}
\uses{lem:geometric-dual-quotient,lem:nonzero-pairings}
\lean{BostConjecture.Adelic.Geometry.exists_canonicalPairing_nonzero_sequence_of_geometricQuotient}
If $q:W\to Q$ is geometrically semistable, there exists $d_0>0$, divisible
by $\dim Q$ and every $a_j$, such that the canonical pairings for
$A_q:Q^\vee\to\bigotimes_jE_j^\vee$ are nonzero over $K$ at every degree
$d_0r$, $r>0$, for suitable rectangular tableaux.
\end{corollary}
\begin{proof}
Apply Lemma~\ref{lem:geometric-dual-quotient}, then
Lemma~\ref{lem:nonzero-pairings} with factors $E_j^\vee$ and source $Q^\vee$.
Finite-dimensional duality preserves dimension, so its divisibility
conditions are exactly those stated here. Its conclusion concerns the
original $K$-linear map $A_q$.
\end{proof}
